\section{Hamiltonian transport and the obstruction}
\label{sec:transport}

We now transport the generating family of
Proposition~\ref{prop:embedded-realization} under a hypothetical Hamiltonian
isotopy.  The construction below keeps track of its actual homogeneous
data at infinity.  Each step adds a fixed quadratic form of split
signature, followed by a homotopy through regular homogeneous data.
In particular, the construction applies to the nonconstant tube family
used here. The construction follows the classical strategy of
composing with small Hamiltonian steps and adding split quadratic
variables \cite[Sections~1.3--1.7]{Sikorav1987}; generating families and their invariance
under isotopy are also developed in \cite{EG1998}.
For a related generating-family lifting statement, see
\mbox{\'{A}lvarez-Gavela} et al.~\cite[Lemma~3.23]{AIS2026}; the argument needed here is given directly.

We use the Morse families and generating maps $\Gamma_F$ of
Section~\ref{sec:embedding}. For any manifold $M$, write $M_0\subset T^*M$
for its zero section. Every family constructed below has
$\Gamma_F$ a diffeomorphism from $\Crit(F)$ onto the Lagrangian it
generates.

\subsection{A Euclidean cylinder and a compactly supported extension}

Retain $N=2m$ and $Q=S^9\times S^{N-1}$ from the preceding construction,
and put $d=10+N$.  View $Q$ as the product of unit spheres in
$\R^{10}\times\R^N=\R^d$.  For $0<\nu<1$, the map
\[
 \iota:Q\times(-\nu,\nu)^2\longrightarrow\R^d,
 \qquad
 \iota((b,v),\alpha)=((1+\alpha_1)b,(1+\alpha_2)v)
\]
is a diffeomorphism onto a bounded open neighborhood $U$ of $Q$.
Its inverse, the projection $r_Q:U\to Q$, and its two radial
coordinates extend smoothly to a neighborhood of $\overline U$.
The cotangent lift gives the canonical identification
\[
 T^*U\cong T^*Q\times T^*((-\nu,\nu)^2),
 \qquad
 \lambda_U=\lambda_Q+p_\alpha\cdot d\alpha.
\]
If $F$ is the family of Proposition~\ref{prop:embedded-realization},
then
\begin{equation}
 F^{(0)}(x,w)=F(r_Q(x),w)
 \label{eq:transport-initial-family}
\end{equation}
is a Morse family whose generated image is the cylinder
\begin{equation}
 \mathcal L_0=
 \{(q,\alpha;p_q,0):(q,p_q)\in L,\ \alpha\in(-\nu,\nu)^2\}.
 \label{eq:transport-cylinder}
\end{equation}
Its critical parametrization is a diffeomorphism.  All its vertical
critical points lie in a uniform ball, and it has the stated smoothness
and homogeneous estimates uniformly up to $\overline U$.

\begin{lemma}[Extension of the reverse isotopy]
\label{lem:transport-cylinder-extension}
Suppose that a smooth compactly supported Hamiltonian isotopy of $T^*Q$
sends $Q_0$ to $L$.  There is a smooth Hamiltonian on $T^*\R^d$,
supported for all times in one compact subset of $T^*U$, whose flow
$\widetilde\varphi_t$ has the following properties.  Every trajectory
starting on $\mathcal L_0$ has constant radial position $\alpha$.
For some $\nu_c>0$, the final image over
$U_c=\iota(Q\times(-\nu_c,\nu_c)^2)$ is precisely the zero section:
\begin{equation}
 \widetilde\varphi_1(\mathcal L_0)\cap T^*U_c=(U_c)_0.
 \label{eq:transport-central-zero}
\end{equation}
\end{lemma}

\begin{proof}
Let $\varphi_t$ be the assumed original flow, generated by $H_t$.
The path $\varphi_{1-t}\circ\varphi_1^{-1}$ starts at the identity,
ends at $\varphi_1^{-1}$, and is generated by
\begin{equation}
 H_t^Q=-H_{1-t}.
 \label{eq:transport-reverse-sign}
\end{equation}
Indeed differentiation of this path gives $-X_{H_{1-t}}$ at its
current point; the fixed composition on the right introduces no
conjugation of the Hamiltonian.

Choose $0\le\zeta\le1$ supported in $(-\nu,\nu)^2$ and identically
$1$ on $(-\nu_c,\nu_c)^2$.  Choose $0\le\xi\le1$ of compact support
in $\R^2$, with $\xi=1$ on a ball to be specified.  On $T^*U$ set
\begin{equation}
 \widetilde H_t(q,\alpha;p_q,p_\alpha)
   =H_t^Q(q,p_q)\zeta(\alpha)\xi(p_\alpha),
 \label{eq:transport-extended-hamiltonian}
\end{equation}
and extend by zero to $T^*\R^d$.  This is smooth with uniform compact
support: $Q$ is compact, $H_t^Q$ has uniformly compact momentum
support, and both additional coordinate pairs are cut off.

The convention $\iota_{X_H}d\lambda=-dH$ gives the radial equations
\begin{equation}
 \dot\alpha=H_t^Q\zeta\,\nabla\xi(p_\alpha),
 \qquad
 \dot p_\alpha=-H_t^Q\nabla\zeta\,\xi(p_\alpha).
 \label{eq:transport-radial-equations}
\end{equation}
Set
\[
 M=\sup_{t,q,p_q}\abs{H_t^Q(q,p_q)}\,\sup_\alpha\norm{\nabla\zeta(\alpha)}.
\]
This number is independent of the size of the plateau of $\xi$.
Take $\xi=1$ on $\{\norm{p_\alpha}\le M+1\}$.  A trajectory starting on
$\mathcal L_0$ starts with $p_\alpha=0$, and the second equation gives
$\norm{p_\alpha(t)}\le Mt$ for $0\le t\le1$.  Hence it stays in this
plateau, the first equation vanishes, and $\alpha$ is constant.
When $\alpha\in(-\nu_c,\nu_c)^2$, also $\nabla\zeta=0$; then
$p_\alpha=0$ and the $T^*Q$ motion is exactly the reverse isotopy.
It sends $L$ to $Q_0$.  Since all radial positions on the entire
cylinder are fixed, no other part of that cylinder enters the central
slab.  This proves \eqref{eq:transport-central-zero}.
\end{proof}

\subsection{A compact mixed generating function for each small step}

Fix a compact subset $K\Subset T^*U$ containing the support of
$\widetilde H_t$ in its interior for every $t$.  Points outside $K$
are stationary, and uniqueness of flow lines prevents a trajectory
from crossing into the stationary exterior in finite time.  In
particular, all the flow maps and their successive time increments
are supported in this common compact region.  Smoothness in time and
compactness allow a finite subdivision
\[
 0=t_0<t_1<\cdots<t_s=1
\]
such that every map
$\psi_j=\widetilde\varphi_{t_j}\widetilde\varphi_{t_{j-1}}^{-1}$
is as close to the identity in $C^1$ as required below.

\begin{lemma}[Mixed coordinates and the action primitive]
\label{lem:transport-mixed-generator}
For a sufficiently small increment, write
$\psi(x,p)=(X,P)$ in the Euclidean cotangent coordinates.  The pair
$(X,p)$ is a global coordinate system on its graph.  There is a smooth
compactly supported function $\sigma(X,p)$ such that
\begin{equation}
 x=X+\sigma_p(X,p),\qquad P=p+\sigma_X(X,p).
 \label{eq:transport-mixed-graph}
\end{equation}
One can choose a single compact $D\Subset U$ containing the spatial
support projection of all these functions, independently of a
subsequent refinement of the subdivision.  Also
$I+\sigma_{pX}$ is invertible, and refinement makes
$\norm{\sigma_p}_\infty$ arbitrarily small.
\end{lemma}

\begin{proof}
Require $\norm{D_xX-I}<\varepsilon<1$ uniformly.  For each fixed $p$,
the map $x\mapsto X(x,p)$ is a local diffeomorphism and satisfies
\[
 \norm{X(x_1,p)-X(x_2,p)}\ge(1-\varepsilon)\norm{x_1-x_2}.
\]
It is injective and proper, since it is the identity outside a compact
set.  Its image is both open and closed in $\R^d$, so it is onto.
The inverse $x=x(X,p)$ is smooth jointly in $(X,p)$.

For a step from time $a$ to time $b$, let $\varphi_{s,a}$ denote
its partial flow.  Cartan's formula and the sign convention give
\[
 \mathcal L_{X_{\widetilde H_s}}\lambda
   =d\bigl(\lambda(X_{\widetilde H_s})-\widetilde H_s\bigr).
\]
Consequently the action integral
\begin{equation}
 A(x,p)=\int_a^b
  \bigl(\lambda(X_{\widetilde H_s})-\widetilde H_s\bigr)
          (\varphi_{s,a}(x,p))\,ds
 \label{eq:transport-action}
\end{equation}
satisfies $\psi^*\lambda-\lambda=dA$.  It is compactly supported:
outside $K$ the trajectory is stationary and its integrand is zero.
On the graph, expressed in mixed coordinates, put
\[
 S(X,p)=A(x(X,p),p)+x(X,p)\cdot p.
\]
Then
\[
 dS=P\cdot dX+x\cdot dp.
\]
Writing $S=X\cdot p+\sigma$ gives
\eqref{eq:transport-mixed-graph}.  Explicitly,
\begin{equation}
 \sigma(X,p)=A(x(X,p),p)+(x(X,p)-X)\cdot p.
 \label{eq:transport-compact-sigma}
\end{equation}
This expression vanishes outside the compact mixed-coordinate image
of $K$.  Its spatial projection lies in a fixed compact subset of
$U$, because both the input and output trajectories remain in $K$.
Enlarge that projection slightly to obtain $D\Subset U$.  Thus for
every $X\in U\setminus D$, $\sigma$ and all its derivatives vanish
on a neighborhood of $X$, for all momenta.  Finally
$I+\sigma_{pX}=D_Xx$ is invertible, and
$\sigma_p=x-X$ tends uniformly to zero with the length of the time
increment.
\end{proof}

\subsection{Transport of one Morse family}

The following step supplies the analytic induction.  The variables
called $w$ include all the stabilizing variables introduced at earlier
steps.

\begin{proposition}[One transport step with controlled infinity data]
\label{prop:transport-step}
Suppose $F_o:U\times\R^r\to\R$ is a smooth Morse family, with
$\Gamma_{F_o}$ a diffeomorphism onto its generated Lagrangian
$\mathcal L_o\subset T^*U$.  Suppose that it is smooth up to
$\overline U$, with uniform bounds on compact vertical sets, and
equals a homogeneous function $g_o(x,w)$ for $\norm{w}\ge R_o$, where
$g_o$ is smooth for $w\ne0$ and
\begin{equation}
 \norm{\nabla_wg_o(x,w)}\ge c_o\norm{w},
 \qquad \norm{d_xg_o(x,w)}\le C_o\norm{w}^2,
 \qquad c_o>0.
 \label{eq:transport-old-bounds}
\end{equation}
For a small step $\psi$ as in
Lemma~\ref{lem:transport-mixed-generator}, there is a smooth Morse
family on $U\times\R^{r+2d}$ generating $\psi(\mathcal L_o)$ with
the same properties.  Its critical locus corresponds
diffeomorphically to $\Crit(F_o)$.  On every fixed compact parameter
slice in $U$, its homogeneous data represent the split quadratic
stabilization of the preceding tube class.
\end{proposition}

\begin{proof}
Introduce $z,p\in\R^d$. Before choosing cutoffs, consider the local expression
\[
 F_o(X+z,w)+\sigma(X,p)-z\cdot p,
 \qquad X+z\in U.
\]
Its $p$- and $z$-equations are $z=\sigma_p(X,p)$ and
$p=\nabla_xF_o(X+z,w)$. Together with the old $w$-equation,
these are precisely the mixed graph equations for composition with
$\psi$: $z$ is the displacement from the new base point $X$ to the
old one, and $p$ is the old momentum. The cutoffs below keep the
shifted base point in $U$ and switch off the shift for large
$\norm{w}$. We will then compare the result with
$g_o(X,w)-z\cdot p$ to restore homogeneity at infinity.

\emph{Choice of cutoffs.}
Choose $\zeta_1:U\to[0,1]$ of compact support, identically $1$ on
an open set $V$ containing $D$.  Fix $\delta_0>0$ so that
\begin{equation}
 2\delta_0<\operatorname{dist}(\operatorname{supp}\zeta_1,\partial U),
 \qquad \norm{\sigma_p}_\infty<\delta_0.
 \label{eq:transport-buffer}
\end{equation}
The first choice is made before refining the subdivision to obtain
the second.  Let $\beta_1:\R^d\to[0,1]$ be $1$ for
$\norm{z}\le\delta_0$ and $0$ for $\norm{z}\ge2\delta_0$.
Choose a smooth radial $\chi:\R^r\to[0,1]$, equal to $1$ on a ball
strictly larger than $\{\norm{w}\le R_o\}$, of compact support, and such
that
\begin{equation}
 \norm{\nabla\chi(w)}\le\delta_1/\norm{w}\quad(w\ne0).
 \label{eq:transport-log-cutoff}
\end{equation}
Here $\delta_1>0$ may be arbitrarily small.  For example, after fixing
the inner radius $R_1>R_o$, take
$\chi(w)=\vartheta(\log(\norm{w}/R_1)/T)$ outside that radius, where
$\vartheta$ is a fixed smooth cutoff equal to $1$ near $(-\infty,0]$
and $0$ near $[1,\infty)$; choosing $T$ sufficiently large gives
\eqref{eq:transport-log-cutoff}.  Extend it by $1$ at smaller radii.

Define
\begin{align}
 \kappa(X,z,w)&=\zeta_1(X)\beta_1(z)\chi(w),
 &x&=X+\kappa(X,z,w)z,\label{eq:transport-kappa}\\
 \widehat F(X,w,z,p)
   &=F_o(X+\kappa(X,z,w)z,w)+\sigma(X,p)-z\cdot p.
 \label{eq:transport-shifted-family}
\end{align}
The shifted argument always lies in $U$.  Indeed, when the shift is
nonzero, $X\in\operatorname{supp}\zeta_1$ and its length is less
than $2\delta_0$; the buffer in \eqref{eq:transport-buffer} places
that entire ball in $U$.  When the shift vanishes the assertion is
immediate.

\emph{The full critical equations.}
Write $a=\nabla_xF_o(x,w)$, identifying Euclidean covectors and
vectors.  Direct differentiation, before discarding any cutoff
terms, gives
\begin{align}
 \widehat F_p&=\sigma_p(X,p)-z,\label{eq:transport-critical-p}\\
 \widehat F_z&=\kappa a+(a\cdot z)\nabla_z\kappa-p,
       \label{eq:transport-critical-z}\\
 \widehat F_w&=\nabla_wF_o(x,w)+(a\cdot z)\nabla_w\kappa,
       \label{eq:transport-critical-w}\\
 \widehat F_X&=a+(a\cdot z)\nabla_X\kappa+\sigma_X(X,p).
       \label{eq:transport-output}
\end{align}
At any vertical critical point, \eqref{eq:transport-critical-p}
forces $z=\sigma_p(X,p)$, hence $\norm{z}<\delta_0$ and $\beta_1=1$
on a neighborhood of that point.  When $\norm{w}\ge R_o$, the old
function is homogeneous at every allowed shifted base point.
Equations \eqref{eq:transport-old-bounds} and
\eqref{eq:transport-log-cutoff} therefore imply
\begin{equation}
 \norm{\widehat F_w}
 \ge(c_o-C_o\delta_0\delta_1)\norm{w}.
 \label{eq:transport-outer-gradient}
\end{equation}
Choose $\delta_1$ so that $C_o\delta_0\delta_1<c_o/2$.
Thus no vertical critical point has $\norm{w}\ge R_o$; at every one of
them $\chi=1$ on a neighborhood as well.

There are two open regions to consider, namely $V$ and $U\setminus D$;
they cover $U$.
For $X\in V$, the three cutoff factors are $1$ near a solution, and
the critical equations become
\begin{equation}
 \nabla_wF_o(x,w)=0,
 \qquad p=\nabla_xF_o(x,w),
 \qquad x=X+\sigma_p(X,p),
 \qquad x=X+z.
 \label{eq:transport-supported-critical}
\end{equation}
They express exactly composition with the graph
\eqref{eq:transport-mixed-graph}, and the generated covector is
\begin{equation}
 \widehat F_X=p+\sigma_X(X,p)=P.
 \label{eq:transport-supported-output}
\end{equation}
The critical equation is regular here.  In local total coordinates
$(x,w,p,X)$, first impose the regular old equation
$\nabla_wF_o=0$.  With $(x,w)$ fixed, the equation
$p-\nabla_xF_o=0$ has identity coefficient in $p$.  With $(x,w,p)$
fixed, the last graph equation has invertible coefficient
$I+\sigma_{pX}$ in $X$.  This triangular order proves surjectivity
of the differential of the full new vertical-gradient map.

For $X\in U\setminus D$, the function $\sigma$ vanishes on a spatial
neighborhood for all $p$.  The equations instead reduce to
\begin{equation}
 z=0,\qquad x=X,\qquad
 \nabla_wF_o(X,w)=0,\qquad
 p=\zeta_1(X)\nabla_xF_o(X,w).
 \label{eq:transport-outside-critical}
\end{equation}
The auxiliary variable $p$ may therefore differ from the old
momentum.  Nevertheless, \eqref{eq:transport-output} gives exactly
\begin{equation}
 \widehat F_X=\nabla_xF_o(X,w),
 \label{eq:transport-outside-output}
\end{equation}
because all the spatial cutoff terms are multiplied by $z=0$.
This is the identity correspondence, which is the true graph of
$\psi$ over this region.  For regularity, first use the equation
$-z=0$ to eliminate $z$, then the equation
$\zeta_1\nabla_xF_o-p=0$ to eliminate $p$; the remaining equation
is the regular old vertical-gradient equation.  This also checks
the entire transition region of $\zeta_1$.  The two descriptions
agree on the overlap, where $\zeta_1=1$.

\emph{The correspondence of critical loci.}
For an old critical point $(x,w)$ put
$p_o=\nabla_xF_o(x,w)$ and $(X,P)=\psi(x,p_o)$.  Define
\begin{equation}
 (x,w)\longmapsto
 \bigl(X,w,z=x-X,\ p=\zeta_1(X)p_o\bigr).
 \label{eq:transport-critical-correspondence}
\end{equation}
If $X\in V$, these are the mixed graph equations, with $p=p_o$.
If $X\notin D$, the true graph is the identity, so $x=X$, $z=0$,
and \eqref{eq:transport-outside-critical} applies.  Thus
\eqref{eq:transport-critical-correspondence} is a smooth map into
$\Crit(\widehat F)$ on all of $\Crit(F_o)$.  Its inverse is
\begin{equation}
 (X,w,z,p)\longmapsto(X+\kappa(X,z,w)z,w).
 \label{eq:transport-critical-inverse}
\end{equation}
The preceding equations show that these maps are inverse and that
the generated image is exactly $\psi(\mathcal L_o)$.
The step preserves $T^*U$, since its support is compactly contained
there.  Thus no old or new critical point is lost at the base
boundary, and $\Gamma_{\widehat F}$ is a diffeomorphism onto the transported
image.  All new solutions obey a uniform bound: $\norm{w}<R_o$,
$\norm{z}<\delta_0$, and
\[
 \norm{p}\le
 \sup\{ \norm{\nabla_xF_o(x,w)}:x\in\overline U, \norm{w}\le R_o\}<\infty.
\]

\emph{Restoring exact homogeneity at infinity.}
Set $W=(w,z,p)$ and consider the homogeneous reference
\begin{equation}
 G_0(X,W)=g_o(X,w)-z\cdot p.
 \label{eq:transport-homogeneous-reference}
\end{equation}
It is $C^1$ everywhere, homogeneous of degree two, and satisfies
\begin{equation}
 \norm{\nabla_WG_0}^2
  =\norm{\nabla_wg_o}^2+\norm{z}^2+\norm{p}^2
  \ge\min(c_o^2,1)\norm{W}^2.
 \label{eq:transport-reference-gradient}
\end{equation}
Both the value error and its \emph{full} vertical first derivative
are uniformly bounded.  To verify this explicitly, write
\begin{align}
 E(X,W)&=\widehat F(X,W)-G_0(X,W)\notag\\
  &=\bigl[F_o(X+\kappa z,w)-F_o(X,w)\bigr]
      +\bigl[F_o(X,w)-g_o(X,w)\bigr]+\sigma(X,p).
 \label{eq:transport-error-decomposition}
\end{align}
The first bracket and every term in its vertical derivative can be
nonzero only with $w,z$ in fixed bounded sets, where the required
derivatives of $F_o$ are uniformly bounded.  The second bracket and
its vertical derivative are bounded because it vanishes for
$\norm{w}\ge R_o$.  The final term has compact support.  The expression
$-z\cdot p$, although unbounded, has cancelled exactly.  Thus for
some finite $M_0,M_1$,
\begin{equation}
 \abs{E(X,W)}\le M_0,
 \qquad \norm{\nabla_WE(X,W)}\le M_1
 \quad(X\in\overline U, W\in\R^{r+2d}).
 \label{eq:transport-bounded-error}
\end{equation}
No compact-support assertion for $E$ in all its variables is needed.

Apply Lemma~\ref{lem:stabilization-smoothing} with
$q'(z,p)=-z\cdot p$, uniformly over $\overline U$. Its homogeneous
smoothing is
\begin{equation}
 G_\eta(X,W)=
   \rho_\eta(\norm{w}/\norm{W})g_o(X,w)-z\cdot p
 \quad(W\ne0),
 \label{eq:transport-angular-smoothing}
\end{equation}
where $\rho_\eta$ vanishes for arguments at most $\eta$ and is $1$
for arguments at least $2\eta$. Choose $\eta$ small and then $A$
sufficiently large as in the lemma, with $A$ beyond the uniform
critical-locus bound. Let $\tau_A$ be the radial cutoff from that
construction and set
\begin{equation}
 F_n=(1-\tau_A)\widehat F+\tau_A G_\eta.
 \label{eq:transport-final-step-family}
\end{equation}
The lemma shows that this function is smooth, equals $G_\eta$ for
$\norm{W}\ge2A$, and has no new critical points. It leaves a
neighborhood of the entire critical locus unchanged, so $F_n$ has
the same regular critical locus and generated parametrization as
$\widehat F$. It also supplies a straight homotopy from $G_0$ to
$G_\eta$ through regular $C^1$ homogeneous functions.

The new homogeneous function $G_\eta$ is smooth off zero, smoothly
dependent on $X$ up to $\overline U$, and has a uniform positive
gradient margin on the unit sphere.  Homogeneity and compactness
there give the next constants of the form
\eqref{eq:transport-old-bounds}; the remaining uniform smooth bounds
follow directly from the construction.  The quadratic form
$-z\cdot p$ is nondegenerate with signature $(d,d)$: in fixed linear
coordinates $z=a+b$, $p=a-b$ it equals $-\norm{a}^2+\norm{b}^2$.
The space of nondegenerate forms of signature $(d,d)$ is connected.
Choose a fixed path in this space from $-z\cdot p$ to the standard
split form in the chosen added coordinates. Its direct sum with
$g_o(X,w)$ has nonzero total gradient off the origin, since both block
gradients can vanish only there. This path is independent of $X$, so
direct sum defines a homotopy over the whole base. Together with the
straight homotopy to $G_\eta$, it preserves the free stable tube class
on every fixed compact parameter slice.
This proves all the assertions.
\end{proof}

\subsection{The final critical point and the contradiction}

\begin{proof}[Proof of Theorem~\ref{thm:main}]
Take the class $a\ne0$ with $\gamma_*a=0$ supplied by
Proposition~\ref{prop:kernel-class}, and the closed exact embedded
Lagrangian $L\subset T^*(S^9\times S^{N-1})$ supplied by
Proposition~\ref{prop:embedded-realization}.  Suppose that a
Hamiltonian isotopy as in the theorem sent the zero section to $L$.
Lemma~\ref{lem:transport-cylinder-extension} supplies a compactly
supported Euclidean Hamiltonian whose final cylinder is the zero
graph over $U_c$.  Start with
\eqref{eq:transport-initial-family} and apply
Proposition~\ref{prop:transport-step} to each of its $s$ increments.
Only finitely many steps are required. The subdivision is fixed by the
mixed-coordinate and spatial-buffer conditions, which depend on the
Hamiltonian and its support. At each step the new family constants
change only the logarithmic cutoff, smoothing parameter and outer
radius; they impose no further restriction on the step length.
Thus the estimates and cutoff radii may be chosen anew at each step.
We obtain a smooth family
\[
 F^{(s)}:U\times\R^{2p}\longrightarrow\R,
 \qquad p=m+ds,
\]
with exactly homogeneous smooth-off-zero data at infinity, whose
regular critical locus maps diffeomorphically onto
$\widetilde\varphi_1(\mathcal L_0)$.

Over $U_c$, the full critical-locus projection
\begin{equation}
 \Crit(F^{(s)})\big|_{U_c}\longrightarrow U_c
 \label{eq:transport-final-projection}
\end{equation}
is a diffeomorphism: it is the generating parametrization followed
by projection of the zero graph.  In particular there is precisely
one critical point over every $X\in U_c$, and these points form a
smooth section.  Their vertical Hessians are invertible.  Indeed,
if $v\in\R^{2p}$ were in the kernel of such a Hessian, then
\[
 D(\partial_WF^{(s)})_{(X,W)}(0,v)
   =\partial_W^2F^{(s)}(X,W)v=0.
\]
Since the full critical equation is regular, its zero set has this
kernel as its tangent space.  Thus $(0,v)$ would be tangent to
$\Crit(F^{(s)})$ and in the kernel of the differential of
\eqref{eq:transport-final-projection}.  That projection is a
diffeomorphism, so $v=0$; the square Hessian is nonsingular.

Fix $v_0\in S^{N-1}$ and restrict the base to
\[
 b\longmapsto\iota((b,v_0),0),\qquad b\in S^9.
\]
The resulting smooth family $f_b:\R^{2p}\to\R$ therefore has
exactly one nondegenerate fiber critical point, varying smoothly
with $b$.  Its homogeneous data belong to $T_{p,p}$, since each
transport step added signature $(d,d)$.  On this fixed slice the
initial homogeneous family was $g_b$, and every subsequent step
preserved its class after split stabilization.  Hence, writing
$g_\infty:S^9\to\mathbf T$ for the final stable tube family,
\[
 [g_\infty]_{\mathrm{free}}=[a]_{\mathrm{free}}\ne[*],
 \qquad \gamma\circ g_\infty\simeq *.
\]
This formulation does not require the smoothing homotopies to fix a
chosen basepoint.  Equivalently, the basepoint paths supplied by
these homotopies identify the resulting based class with $a$.
All hypotheses of Proposition~\ref{prop:one-critical} hold, including
the uniform homogeneous behavior and the required large indices.
That proposition makes the same tube family freely null in the
stable tube space, contradicting the non-null class of
Proposition~\ref{prop:kernel-class}.  The assumed compactly supported
Hamiltonian isotopy does not exist.  This closed exact embedded $L$
is therefore the required counterexample.
\end{proof}
